\documentclass[11pt]{article}

\usepackage[margin=1.1in]{geometry}
\usepackage{amsmath,amssymb,amsthm}
\usepackage{stmaryrd}
\usepackage{mathtools}
\usepackage{booktabs}
\usepackage{array}
\usepackage{xcolor}
\usepackage[colorlinks=true,linkcolor=black,citecolor=black,urlcolor=blue]{hyperref}
\usepackage[expansion=false]{microtype}

\theoremstyle{plain}
\newtheorem{theorem}{Theorem}
\newtheorem{proposition}[theorem]{Proposition}
\newtheorem{lemma}[theorem]{Lemma}
\newtheorem{corollary}[theorem]{Corollary}
\newtheorem{conjecture}[theorem]{Conjecture}
\theoremstyle{definition}
\newtheorem{definition}[theorem]{Definition}
\newtheorem{remark}[theorem]{Remark}
\newtheorem{example}[theorem]{Example}
\newtheorem{hypothesis}[theorem]{Hypothesis}

% --- typography: the rho calculus is NEVER written with the Greek letter ---
\newcommand{\rhoc}{\textsf{rho}}
\newcommand{\GSLT}{\mathbf{GSLT}}
\newcommand{\Kc}{\mathsf{K}}
\newcommand{\Gg}{\mathbb{G}}
\newcommand{\Lo}{\mathcal{L}}
\newcommand{\Obs}{\mathrm{Obs}}
\newcommand{\Ask}{\mathsf{Ask}}
\newcommand{\Get}{\mathsf{Get}}
\newcommand{\Bld}{\mathsf{Build}}
\newcommand{\args}{\mathsf{args}}
\newcommand{\bis}{\sim}
\newcommand{\Terms}{\mathcal{T}}
\newcommand{\Ob}{\mathcal{O}}

\title{\bfseries Every Constructor an Interaction\\[.35em]
\large Observer extensions, and the discriminating power\\
of spatial formulae in interactive graph-structured lambda theories}
\author{L.\ G.\ Meredith\\ \small F1R3FLY.io}
\date{August 7, 2026 --- second draft}

\begin{document}
\maketitle

\begin{abstract}
Spatial logics separate more than bisimulation does. This is not a defect and it is not news: Caires
observed two decades ago that the equivalence induced by a spatial logic for the $\pi$-calculus sits
strictly between structural congruence and strong bisimilarity, and the intensionality of such logics
was settled shortly afterwards by Sangiorgi and by Hirschkoff, Lozes and Sangiorgi. Yet the same
observation, restated, is routinely taken to refute any claim that a generated spatial--behavioural
logic is \emph{adequate} --- that its formulae separate exactly what interaction separates. The two
claims look incompatible and are not.

They are not incompatible because there is no single observational bisimulation against which a
spatial logic must be measured. Observational equivalence is indexed by the destructuring capabilities
admitted to the observer. We make the index explicit and mechanical. Given an interactive
graph-structured lambda theory $\Gg$ with interaction constructor $\Kc$, and a set $\Ob$ of
constructors the observer may take apart, we adjoin one nullary probe atom per constructor and per
argument position, freely adjoin a correctly sorted argument former per signature profile, and overload
$\Kc$ so that a probe meets a term. Three rewrite rules per constructor then make asking, projecting
and rebuilding into interactions. In the resulting theory $\Obs(\Gg,\Ob)$ the observer's structural
questions are ordinary transitions, and bisimilarity descends monotonically from context-labelled
bisimilarity on $\Gg$ at $\Ob = \varnothing$ to the equational theory of $\Gg$ at $\Ob = \Sigma$,
which is where the spatial connectives already were. The two disputed propositions are the two
endpoints of one family.

The phenomenon is folklore in the spatial logics community, where extensionality of spatial
observation has been established \emph{semantically}, by exhibiting models --- typically distributed
calculi with failures --- whose observational equivalence is characterised by a spatial logic. Our
handling is \emph{syntactic}: a transformation of presentations, mechanical in the signature. We prove
a reconstruction theorem for the operational side and a separation theorem for the logical side, keeping
them apart, so that a weakness in one does not look like a weakness in both; and we record what the
construction costs. It is not object-capability safe, since unrestricted opening destroys
encapsulation. We describe the shape of the capability discipline that repairs this, and do not develop
it here.

The note is downstream of \cite{BIHOL}, in preparation, which supplies the label discipline: labels are
minimal contexts enabling a rewrite, in the sense of Leifer and Milner, in an open form that does not
restrict to ground terms.
\end{abstract}

\tableofcontents

\section{Introduction}
\label{sec:intro}

\subsection{The phenomenon}

A spatial--behavioural logic for a process calculus carries two layers. The behavioural layer is
modal: formulae speak of the steps a term can take. The structural, or spatial, layer speaks of how the
term is put together --- that it is a parallel composition of two parts with such-and-such properties,
that it is a send on a name, that it is empty.

The layers do not have the same separating power, and the direction of the gap surprises people the
first time they meet it. Behavioural equivalence is coarse by design: it identifies terms that no
experiment can tell apart. Structural predicates are finer, because they can ask questions no
experiment can pose --- or rather, no experiment \emph{in the calculus as presented}. That
qualification is the subject of this note.

\subsection{The apparent contradiction}

When the logic is not designed but \emph{generated} from the presentation of the calculus --- one
structural connective per term former, one modality per rewrite rule and redex position --- one wants
an adequacy theorem: logical equivalence coincides with the calculus's observational equivalence.
Adequacy is what makes a generated logic worth generating. Too coarse and there are observable
distinctions the logic cannot state; too fine and the logic invites its user to chase differences no
experiment could settle.

So one wants to assert both of:
\begin{align}
\text{logical equivalence} \;&=\; \text{observational equivalence}, \tag{A} \\
\text{the spatial layer} \;&\subsetneq\; \text{observational equivalence}. \tag{B}
\end{align}
Read carelessly these are inconsistent, and a careful reader will say so. We have seen (B) offered as
a refutation of (A) more than once, by readers who knew the literature well enough to recognise (B) and
not well enough to know what answers it.

\subsection{The resolution}

There is no such thing as \emph{the} observational equivalence of a calculus independent of what the
observer is permitted to do. (B) is a statement about bisimilarity in $\Gg$, where the observer may
build contexts out of $\Gg$'s constructors and watch what happens. (A) is a statement about
bisimilarity in a theory in which the observer additionally holds instruments for taking terms apart.
Exhibit the index and the family interpolates between them; there is then nothing to reconcile.

The construction is the subject of \S\ref{sec:construction}. Where the cut has a transparent unit its
slogan is:

\begin{quote}
\itshape Make every term former an interaction. Bisimulation observes interactions. Therefore
bisimulation observes every term former.
\end{quote}

\subsection{What is new here, and what is not}
\label{sec:contribution}

The phenomenon is known. \S\ref{sec:lineages} sets out four lineages of prior work, of which one ---
extensionality of spatial observation --- already establishes that spatial logics need not be counted
intensional, and another --- reactive systems --- supplies the label discipline we use. We claim no
priority over any of it, and the underlying idea of reifying constructors as tags with observation
actions is elementary enough that we expect precedents we have not found.

What we add is a change of method, and the change buys generality. The prior extensionality results are
\emph{semantic}: they exhibit a calculus, typically a distributed one with local and remote
communication and partial failures, and show that its observational equivalence is characterised by a
spatial logic. The extra separating power comes from features of that calculus, which are modelling
choices about a particular domain.

Our construction is \emph{syntactic}. It is a transformation of presentations. It reads the signature,
adjoins probe atoms and argument formers, adds three rules per constructor, and stops. It is defined
for any interactive GSLT satisfying the hypotheses of \S\ref{sec:prelim} rather than for one calculus,
and because it is mechanical in the presentation it composes with the machinery that generates the
logic from the same presentation. That last point is what matters for our purposes: if both the logic
and the enlargement are functions of the signature and of $\Ob$, then adequacy is the statement that
two functions of the same argument agree, rather than a coincidence to be checked calculus by calculus.

\subsection{Three things kept apart}

An earlier draft of this note braided three claims together tightly enough that a weakness in one
looked like a weakness in all three. We separate them here and recommend the separation to anyone
reasoning in this area:

\begin{itemize}
\item \textbf{operational reconstruction} --- what the enlarged transition system can distinguish
(\S\ref{sec:recon});
\item \textbf{logical expressiveness} --- what the generated structural language can distinguish
(\S\ref{sec:separation});
\item \textbf{observer capability} --- which instruments a given observer actually holds
(\S\ref{sec:dial}, \S\ref{sec:security}).
\end{itemize}

\subsection{What this note does not claim}

It does not claim that structural inspection was interaction all along. That would be slippery, and an
earlier draft said it. The form of observation remains transition observation, but the set of
observable transitions has been enlarged, and the enlargement is exactly what supplies the additional
discriminating power. The honest statement is that structural observation \emph{can be represented as}
interaction, after conservatively extending the observer's transition theory with destructuring
actions.

It does not develop the capability discipline of \S\ref{sec:security}. It does not treat weak
equivalences or quantitative refinements. And it makes no claim that unrestricted opening is a sensible
thing to permit in a deployed system.

\section{Four lineages}
\label{sec:lineages}

Much of the confusion we are addressing comes from reading distinct strands of prior work as one.

\paragraph{The observation.} Caires, studying an expressive spatial logic for the synchronous
$\pi$-calculus with recursion, gives coinductive and equational characterisations of the equivalence
the logic induces and locates it: strictly between structural congruence and strong bisimilarity
\cite{CairesObs}. The companion development with Cardelli \cite{CairesCardelliI,CairesCardelliII} is
explicit that formulae are taken to denote processes up to structural equivalence rather than up to
bisimulation, precisely because a tensor operator that distinguishes bisimilar processes is difficult
to reconcile with the latter. This is (B), stated by the people who first had to confront it.

\paragraph{The intensionality results.} Sangiorgi \cite{SangiorgiExtInt} and then Hirschkoff, Lozes and
Sangiorgi \cite{HLS02,HLS03} settle where the induced equivalence sits: for logics of this kind it is
essentially structural congruence. Lozes adds the elimination results --- that in the static ambient
logic the adjuncts contribute no expressive power \cite{LozesAdjunct}, later strengthened by Dawar,
Gardner and Ghelli through model-comparison games \cite{DGG} --- and the picture is completed by Caires
and Lozes \cite{CairesLozes}, who show that a minimal dynamic spatial logic already encodes action
quantification and action modalities, so that the spatial and modal fragments have the same separating
power, coinciding with structural congruence modulo permutation of actions. Lozes' thesis
\cite{LozesThesis} is the systematic account. This strand says the logic is \emph{intensional}: it sees
syntax. It is often read, wrongly, as saying that nothing can be done about this.

\paragraph{The extensionality results.} Something can be done about it. Caires and Vieira
\cite{CairesVieira} address the tension directly, showing that there are natural models in which
extensional observational equivalences \emph{are} characterised by spatial logics, composition and void
included --- and concluding that spatial observations need not always be counted intensional, even when
expressive enough to describe structure. Their development is carried out in a minimalist distributed
calculus with local synchronous communication, local computation, asynchronous remote communication,
and partial failures. The additional separating power available to an observer in that setting is what
closes the gap. A related model-theoretic treatment is given by Tuosto and Vieira
\cite{TuostoVieira}; Caires' tutorial survey \cite{CairesSurvey} is the best single entry point.

\paragraph{Reactive systems.} Labelled transition systems were for a long time built by hand, the
researcher choosing labels that would give the right view of behaviour. Leifer and Milner
\cite{LeiferMilner} showed how to derive them: labels should be the \emph{minimal contexts} enabling a
rewrite, minimality determined by a universal property, and bisimilarity over the derived transition
system is then a congruence by construction rather than by contextual closure after the fact. Sewell
\cite{Sewell} and, in a higher-order setting, Di Gianantonio, Honsell and Lenisa \cite{LambdaLTS}
develop the programme; bigraphs \cite{Bigraphs} are its best-known application. This is the strand
that supplies our label discipline, and \cite{BIHOL} is the form in which we use it.

\begin{table}[t]
\centering
\begin{tabular}{@{}p{3.0cm}p{5.3cm}p{5.2cm}@{}}
\toprule
\textbf{strand} & \textbf{question settled} & \textbf{what it does not settle} \\
\midrule
the observation & where the induced equivalence sits relative to $\bis$ and $=_E$ &
whether the gap is intrinsic \\
intensionality & that the induced equivalence is essentially $=_E$; which connectives are eliminable &
whether an observer could be given the missing power \\
extensionality & that in suitable models observational equivalence is characterised spatially &
whether this is obtainable uniformly from a presentation \\
reactive systems & how to derive labels making bisimilarity a congruence &
what happens when the observer is given destructors \\
\bottomrule
\end{tabular}
\caption{The four strands, and the question each leaves open. This note addresses the right-hand cells
of the last two rows.}
\label{tab:lineages}
\end{table}

\begin{remark}[Which bisimilarity]
\label{rem:whichbisim}
The equivalence Caires locates is induced on an action-labelled transition system for a particular
calculus. The equivalence this note works with is the context-labelled one of \S\ref{sec:prelim}. They
are not the same relation, and we do not borrow the earlier statement as though it were ours: claim (B)
is re-established for our relation in \S\ref{sec:nonvacuity}. Conflating them would be exactly the sort
of suppressed index this note exists to expose.
\end{remark}

\section{Preliminaries}
\label{sec:prelim}

We recall what is used and nothing more. For graph-structured lambda theories see \cite{Omnibus}; for
the label discipline, \cite{BIHOL}, which is in preparation and on which this note is frankly
downstream.

\begin{definition}[GSLT]
A \emph{graph-structured lambda theory} is a triple $\Gg = (\Sigma, E, R)$ where $\Sigma$ is the
signature of a lambda theory --- so it may be many-sorted, constructors may bind, and terms may contain
variables --- $E$ is a set of equations, and $R$ is a set of rewrite rules, the graph of rewrites being
adjoined to the structure rather than the structure being enriched over graphs. We write
$\Terms(\Sigma)$ for the terms and $=_E$ for the least congruence containing $E$.
\end{definition}

\begin{definition}[Interactive]
$\Gg$ is \emph{interactive} when $\Sigma$ carries a distinguished binary constructor $\Kc$, the
\emph{interaction constructor} or \emph{cut}, such that the left-hand side of every rule in $R$ is of
the form $\Kc(l_1,l_2)$ up to $=_E$.
\end{definition}

Parallel composition plays this role for CCS, the $\pi$-calculus, the ambient calculus and the \rhoc\
calculus; application plays it for the lambda calculus.

\subsection{Labels are minimal enabling contexts}

Following \cite{LeiferMilner} in the open form of \cite{BIHOL}: a transition $t \xrightarrow{c} d[r]$
is asserted when $l \to r$ is a rule and $d[l] = c[t]$ is an idempotent relative pushout --- so that
$c$ is a \emph{least} context enabling the rewrite, minimality being a universal property in the
relevant slice category. Leifer and Milner's reactive systems fix a distinguished object as the domain
of the IPOs, effectively restricting to ground terms; \cite{BIHOL} observes that the proofs do not
depend on it and adopts an open version, which is what we need, since a lambda theory's terms carry
variables.

Minimality is not a technical convenience. It is what keeps a label from carrying the whole term. A
transition system labelled by \emph{arbitrary} enabling contexts would be so informative that matching
labels would nearly force syntactic agreement, and the equivalence it induced would already be at or
near $=_E$ --- collapsing the distinction this note exists to draw.

\begin{hypothesis}[Redex RPOs]
\label{hyp:rpo}
$\Gg$ has all redex-RPOs, so that the derived transition system exists.
\end{hypothesis}

We inherit this from \cite{BIHOL}, where establishing it for the theories of interest is work in
progress. Everything below is conditional on it.

\subsection{Matching labels}

In a higher-order setting a minimal enabling context still carries \emph{processes} --- for the \rhoc\
calculus the labels are of the form $-\mid \mathrm{in}(n,\lambda x.q)$ with $q$ an arbitrary process.
Matching such labels syntactically yields a relation known to be too fine \cite{SangiorgiHO}; the
standard remedy is to compare the payload up to the bisimilarity being defined.

\begin{definition}[Context-labelled bisimulation]
\label{def:bisim}
A symmetric relation $\mathcal{R}$ is a \emph{context-labelled bisimulation} when $P \mathrel{\mathcal{R}} Q$
implies: for every transition $P \xrightarrow{c} P'$ there are $Q'$ and a label $c'$ with
$Q \xrightarrow{c'} Q'$, $P' \mathrel{\mathcal{R}} Q'$, and $c \mathrel{\mathcal{R}^\ast} c'$, where
$\mathcal{R}^\ast$ compares labels position by position, relating the constructor skeletons by identity
and the process payloads by $\mathcal{R}$. Bisimilarity $\bis_\Gg$ is the greatest such relation.
\end{definition}

\begin{lemma}[Well-definedness]
\label{lem:monotone}
The functional whose greatest fixed point is Definition~\ref{def:bisim} is monotone, so $\bis_\Gg$
exists.
\end{lemma}

\begin{proof}
$\mathcal{R}$ occurs only positively: in the conclusion of the matching clause, and in
$\mathcal{R}^\ast$, which is built from $\mathcal{R}$ by identity and componentwise conjunction. A
monotone functional on the complete lattice of symmetric relations has a greatest fixed point by
Knaster--Tarski.
\end{proof}

\begin{hypothesis}[Congruence]
\label{hyp:cong}
A presentation is \emph{admissible} when it comes equipped with a class $\mathcal{A}$ of contexts,
containing those built from the rewrite constructors, such that $\bis_\Gg$ is a congruence with respect
to $\mathcal{A}$, and observations are restricted to $\mathcal{A}$. We consider only admissible
presentations.
\end{hypothesis}

This is a hypothesis, not a theorem, and deliberately so. We are interested only in context-labelled
transition graphs that make bisimulation a congruence; barbed bisimulation, which obtains congruence by
quantifying over contexts and a calculus-specific notion of observable, does not generalise across
interactive GSLTs and we do not use it. Deriving labels from contexts so that congruence holds by
construction is precisely the point of \cite{LeiferMilner}.

\begin{remark}[The restriction is real]
\label{rem:areal}
$\mathcal{A}$ is not all contexts. For the \rhoc\ calculus, \cite{BIHOL} proves congruence under
$\mathrm{out}(n,-)$, $\mathrm{in}(n,-)$, $-\mid-$ and $\ast(@(-))$, and must exclude contexts with a
hole \emph{under} the quote, such as $\mathrm{out}(@(-),z)$: communication turns on name equality, not
on bisimilarity, so $p \bis q$ with $p \neq q$ gives $@p \neq @q$ and the two contexts behave
differently. \S\ref{sec:reflection} returns to what this costs the present construction.
\end{remark}

\subsection{The generated logic}

\begin{definition}[The generated logic]
$\Lo(\Gg)$ has formulae
\[
\varphi \;::=\; \top \;\mid\; \neg\varphi \;\mid\; \varphi \wedge \psi
\;\mid\; C(\varphi_1,\dots,\varphi_n) \;\mid\; \langle c \rangle \varphi
\]
with one \emph{structural} connective per $n$-ary constructor $C \in \Sigma$ and one \emph{behavioural}
modality per label. Satisfaction for the structural connectives is
\[
P \models C(\varphi_1,\dots,\varphi_n)
\iff
P =_E C(t_1,\dots,t_n) \text{ with } t_i \models \varphi_i \text{ for each } i,
\]
and $P \models \langle c\rangle\varphi$ iff $P \xrightarrow{c} P'$ with $P' \models \varphi$. We write
$\Lo_{\mathrm{s}}(\Gg)$ for the structural fragment and $=_{\Lo}$ for logical equivalence.
\end{definition}

Because $\Kc$ is typically subject to associativity, commutativity and a unit in $E$, the connective
$\Kc(\varphi,\psi)$ is a separating conjunction: it holds of $P$ exactly when $P$ admits \emph{some}
decomposition into two parts satisfying $\varphi$ and $\psi$. The existential over decompositions is the
whole difficulty. It is not a step, so no modality expresses it, and that is where (B) comes from.

\subsection{Admissible equational theories}

The results below need $E$ to admit a well-founded notion of decomposition. We isolate the conditions
rather than discover them mid-proof.

\begin{definition}[Measure]
For $P \in \Terms(\Sigma)$ let $|P|$ be the number of constructor occurrences and
$\mu(P) = \min\{\,|Q| : Q =_E P\,\}$. A representative attaining the minimum is \emph{lean}.
\end{definition}

\begin{definition}[Admissible $E$]
\label{def:adm}
$E$ is \emph{admissible} when:
\begin{description}
\item[(E1) Well-founded.] Every $=_E$-class of closed terms contains a lean representative.
\item[(E2) Hereditary.] If $C(\vec t\,)$ is lean then each $t_i$ is lean, and for $n \ge 1$ we have
$\mu(t_i) < \mu(C(\vec t\,))$.
\item[(E3) Sorted.] $E$ relates only terms of the same sort.
\end{description}
\end{definition}

\begin{remark}[Why (E2) is not free]
\label{rem:unit}
An earlier draft asserted that $E$ is size-respecting on the presentations of interest. It is not: the
unit law $P \mid \mathbf{0} = P$ is not size-respecting, and it is one of the equations doing the most
work below, since it is what makes $P \mid \mathbf{0}$ one of the decompositions the separating
conjunction quantifies over. (E1)--(E2) are the repair. They say nothing about how many decompositions
a term has --- possibly infinitely many --- only that one may always descend along a lean one.
Associativity, commutativity and unit for $\Kc$ satisfy them.
\end{remark}

\section{The observer extension}
\label{sec:construction}

Fix an admissible interactive GSLT $\Gg = (\Sigma,E,R)$ with cut $\Kc$, satisfying
Hypotheses~\ref{hyp:rpo} and \ref{hyp:cong}, with $E$ admissible in the sense of
Definition~\ref{def:adm}.

\subsection{Argument formers, freely adjoined}

$\Kc$ is binary; constructors have arbitrary arity and heterogeneous argument sorts. To present the
arguments of an $n$-ary constructor to $\Kc$ they must be gathered into one operand.

\begin{remark}[Not lists, and not Turing completeness]
\label{rem:noturing}
An earlier draft encoded the arguments as a list and appealed to Turing completeness to supply the list
former. That was wrong twice over. Computability guarantees the ability to represent data; it does not
supply an internal coding function injective modulo an arbitrary $E$, and injectivity modulo $=_E$ is
exactly what the reconstruction below needs. And a general lambda theory is many-sorted, so a
homogeneous list cannot hold heterogeneous arguments without a universal code sort, which is a further
construction choice.

The repair is simpler than the thing repaired: \emph{freely adjoin} an argument former per signature
profile. Free adjunction gives correct sorting and injectivity by construction, introduces no encoding
choice, and removes Turing completeness from the note entirely. Nothing below uses it.
\end{remark}

\begin{definition}[Observer extension]
\label{def:ext}
Let $\Ob \subseteq \Sigma$ be an \emph{observation signature}. The \emph{observer extension}
$\Obs(\Gg,\Ob) = (\Sigma^\Ob, E, R^\Ob)$ is given by:
\begin{itemize}
\item for each $C \in \Ob$ of profile $\tau_1 \times \cdots \times \tau_n \to \tau$, a fresh sort
$A_C$ and a freely adjoined former $\args_C : \tau_1 \times \cdots \times \tau_n \to A_C$;
\item nullary \emph{probe atoms} $\Ask_C$, $\Bld_C$ and $\Get_{C,i}$ for $1 \le i \le n$;
\item $\Kc$ overloaded so that $\Kc(\Ask_C, t)$, $\Kc(\Get_{C,i}, a)$ and $\Kc(\Bld_C, a)$ are well
formed for $t : \tau$ and $a : A_C$;
\item $E$ unchanged;
\item $R^\Ob = R \cup \{\, \mathrm{ask}_C,\ \mathrm{get}_{C,i},\ \mathrm{bld}_C : C \in \Ob \,\}$ where
\begin{align*}
\mathrm{ask}_C   &: \;\; \Kc(\Ask_C,\, C(x_1,\dots,x_n)) \;\longrightarrow\; \args_C(x_1,\dots,x_n),\\
\mathrm{get}_{C,i} &: \;\; \Kc(\Get_{C,i},\, \args_C(x_1,\dots,x_n)) \;\longrightarrow\; x_i,\\
\mathrm{bld}_C   &: \;\; \Kc(\Bld_C,\, \args_C(x_1,\dots,x_n)) \;\longrightarrow\; C(x_1,\dots,x_n).
\end{align*}
\end{itemize}
Rules of $R$ are \emph{proper}; the rest are \emph{administrative}. We write $\Gg^+$ for
$\Obs(\Gg,\Sigma)$.
\end{definition}

Four comments on the shape of this definition.

First, the overloading of $\Kc$ rather than the introduction of a fresh eliminator is the point of the
whole construction. $\Kc$ is not an arbitrary constructor; it is the one the transition relation is
indexed on, the one that defines what it means for two things to meet. A probe placed in interaction
position turns a structural question into a rendezvous, and a rendezvous is what a derived transition
system is built to observe.

Second, every administrative rule is \emph{interaction-headed by construction}: its left-hand side is
literally a $\Kc$. An earlier draft put the constructor being opened on the left of the arrow, which
forced an appeal to a unit for $\Kc$ and a special case, patched by an extra wrapper, for cuts without
one --- and the patch destroyed the very lemma the proof depended on. Putting the probe on the left
removes the case distinction: the observer supplies the context, so no unit is needed anywhere.

Third, the probe atoms are indexed. It is $\Ask_C$ rather than a single $\Ask$ that makes the label of
an opening transition record \emph{which} constructor was found, and $\Get_{C,i}$ rather than a single
$\Get$ that makes projection record \emph{which} argument was taken. This is what carries the
information the reconstruction needs; see Remark~\ref{rem:labelnottarget}.

Fourth, $\mathrm{bld}_C$ is not used in any proof below. It is present so that the extension supplies
construction as well as deconstruction, and so that the observer's instruments form a coherent kit
rather than a one-way mirror.

\begin{proposition}[$\Obs(\Gg,\Ob)$ is an interactive GSLT]
\label{prop:still}
$\Obs(\Gg,\Ob)$ is a GSLT, is interactive with the same cut $\Kc$, and reflects proper transitions: for
$P,P' \in \Terms(\Sigma)$ and $\rho \in R$, $P \xrightarrow{c} P'$ in $\Gg$ iff the same holds in
$\Obs(\Gg,\Ob)$.
\end{proposition}

\begin{proof}
Every rule of $R^\Ob$ has a left-hand side of the form $\Kc(-,-)$: those of $R$ by hypothesis, the
administrative ones by construction. Reflection of proper transitions holds because $R^\Ob \supseteq R$
and the adjoined formers are mentioned in no rule of $R$, so no $R$-redex is created or destroyed by
their presence, and the RPOs computed for $R$-redexes are unchanged.
\end{proof}

\begin{remark}[Observation as a separate label class]
Nothing forces the administrative rules to be treated as ordinary rules. One may instead take them to
constitute a distinct class of labels, matched separately. Because Definition~\ref{def:ext} makes them
interaction-headed anyway, the two treatments agree here, and we take the simpler one. The distinction
would matter for a version of the construction in which opening did not require a probe.
\end{remark}

\begin{remark}[Equations as rewrite pairs]
\label{rem:pairs}
One may prefer a presentation with no equations, orienting each equation of $E$ in both directions and
adding both to the administrative rules. The two presentations agree up to administrative steps. We
keep $E$ as equations because doing so localises the source of nondeterminism: with $E$ retained, the
branching of $\mathrm{ask}_C$ comes from $E$-matching, which is exactly the existential over
decompositions that the separating conjunction quantifies. It is worth stressing that these must be
\emph{pairs}. Orienting the equations in one direction only would let bisimilarity fall below $=_E$ onto
raw syntax, which overshoots: the structural connectives respect $=_E$, so a relation finer than $=_E$
is finer than the logic and adequacy fails from the other side.
\end{remark}

\section{Reconstruction}
\label{sec:recon}

\subsection{Exposure}

\begin{lemma}[Exposure]
\label{lem:exposure}
Let $P \in \Terms(\Sigma^\Ob)$ and $C \in \Ob$ be $n$-ary. Then
\[
P \xrightarrow{\;\Kc(\Ask_C,\,-)\;} \args_C(t_1,\dots,t_n)
\quad\text{iff}\quad
P =_E C(t_1,\dots,t_n),
\]
and $\Kc(\Ask_C,-)$ is a least context enabling $\mathrm{ask}_C$ at $P$. Consequently the
$\Ask_C$-labelled transitions of $P$ enumerate exactly the top-level $E$-decompositions of $P$ at $C$.
Likewise $\args_C(\vec t\,)$ has the single transition $\xrightarrow{\Kc(\Get_{C,i},-)} t_i$ with that
label.
\end{lemma}

\begin{proof}
The redex $\Kc(\Ask_C, C(\vec x))$ is formed from $P$ by the context $\Kc(\Ask_C,-)$ exactly when
$P =_E C(\vec t)$; no smaller context enables it, since the rule's left-hand side requires the probe,
and no larger one is minimal. Matching is up to $=_E$, which is where the enumeration comes from. For
$\mathrm{get}$, $\args_C$ is freely adjoined, so $\args_C(\vec t) =_E \args_C(\vec s)$ iff
$t_i =_E s_i$ for all $i$, and the projection is deterministic.
\end{proof}

This is the lemma that does the work, and it is worth pausing on what it says in the case that
motivated the difficulty. Take $C = \Kc$ with $E$ containing associativity, commutativity and unit.
Then $P$'s $\Ask_\Kc$-labelled transitions are in bijection with the ways of writing $P$ as $Q \mid R$
--- including $P \mid \mathbf{0}$ and $\mathbf{0} \mid P$, and every rebracketing and permutation of a
parallel product. The existential over splits in $P \models \varphi \mid \psi$ has become an existential
over transitions. Everything else is bookkeeping.

\begin{remark}[The label, not the target]
\label{rem:labelnottarget}
An earlier draft said that the atom appearing ``in the label'' records the constructor. It did not: the
atom was in the target, and only the rule name distinguished the constructors. With the probe on the
left of the cut the claim is now literally true --- $\Ask_C$ occurs in the context $\Kc(\Ask_C,-)$,
which is the label --- and the proof below may use it directly.
\end{remark}

\begin{lemma}[Calibration]
\label{lem:calib}
If $P =_E Q$ then $P \bis_{\Obs(\Gg,\Ob)} Q$.
\end{lemma}

\begin{proof}
Transitions are defined up to $=_E$, so $=_E$-equal terms have identical labelled transition sets, and
the identity relation on $=_E$-classes is a bisimulation.
\end{proof}

Trivial as stated, and stated anyway, because it fixes the \emph{lower} bound of the calibration: the
enlargement must not push bisimilarity below the equational theory, or it overshoots the logic. It is
also the property the alternative presentation of Remark~\ref{rem:pairs} must be checked against.

\subsection{The theorem}

\begin{theorem}[Reconstruction]
\label{thm:main}
Let $\Gg$ be admissible with $E$ admissible, and let $\Ob = \Sigma$. Then for closed
$P,Q \in \Terms(\Sigma)$,
\[
P \bis_{\Gg^+} Q \iff P =_E Q .
\]
\end{theorem}

\begin{proof}
($\Leftarrow$) is Lemma~\ref{lem:calib}.

($\Rightarrow$) By strong induction on $\mu(P)$, with induction hypothesis: for all closed $P', Q'$
with $\mu(P') < \mu(P)$, if $P' \bis_{\Gg^+} Q'$ then $P' =_E Q'$. Note that the hypothesis constrains
only the left argument, which is what allows the right-hand witnesses produced below to be arbitrary
representatives rather than lean ones.

Let $P \bis_{\Gg^+} Q$ and choose a lean representative $C(t_1,\dots,t_n)$ of $P$, which exists by
(E1). By Lemma~\ref{lem:exposure},
\[
P \xrightarrow{\;\Kc(\Ask_C,-)\;} \args_C(t_1,\dots,t_n).
\]
Since $P \bis_{\Gg^+} Q$, the term $Q$ must match this step with a label related to
$\Kc(\Ask_C,-)$. That label has no process payload --- $\Ask_C$ is a nullary atom and the hole is the
only other position --- so label matching is identity on the constructor skeleton, and $Q$'s matching
label is $\Kc(\Ask_C,-)$ itself. Distinct constructors carry distinct atoms, so no other opening rule
could have matched. By Lemma~\ref{lem:exposure} again, $Q =_E C(s_1,\dots,s_n)$ for some $\vec s$, and
\[
\args_C(t_1,\dots,t_n) \;\bis_{\Gg^+}\; \args_C(s_1,\dots,s_n).
\]
Now apply $\Get_{C,i}$. By Lemma~\ref{lem:exposure} each side has exactly one transition with label
$\Kc(\Get_{C,i},-)$, to $t_i$ and $s_i$ respectively, so $t_i \bis_{\Gg^+} s_i$ for each $i$.

By (E2) each $t_i$ is lean and, for $n \ge 1$, $\mu(t_i) < \mu(P)$. The induction hypothesis gives
$t_i =_E s_i$, whence $P =_E C(\vec t) =_E C(\vec s) =_E Q$. For $n = 0$ the constructor is nullary and
$P =_E C =_E Q$ directly.
\end{proof}

\begin{remark}[Domains]
Theorem~\ref{thm:main} is a statement about $\bis_{\Gg^+}$ restricted to closed $\Sigma$-terms. The
relation itself is defined on all of $\Terms(\Sigma^{\Sigma})$, where it also identifies, for instance,
$\args_C(\vec t)$ with $\args_C(\vec s)$ exactly when the components agree. We are careful about this
in \S\ref{sec:idem}, where an earlier draft was not.
\end{remark}

\begin{remark}[Infinite branching]
A term with replication has infinitely many $E$-decompositions and therefore infinitely many
$\Ask_\Kc$-labelled transitions. Nothing in the proof requires image-finiteness: the induction is on
$\mu$, not on transition depth, and no Hennessy--Milner style characterisation is invoked.
\S\ref{sec:separation} is what makes this safe on the logical side as well.
\end{remark}

\subsection{Binding}

\begin{remark}[Scope]
\label{rem:binding}
For a binding constructor the arguments are abstractions rather than closed terms, and
Theorem~\ref{thm:main} as proved does not cover them: it would need $\alpha$-equivalence to be part of
$E$ and (E1)--(E2) to hold of the resulting classes, a notion of transition beneath a binder, and a
statement of what $\Kc(\Ask_C,-)$ means when $C$ binds. We expect the lambda-theoretic setting to
supply all three --- an abstraction is a term of the theory, so the induction has somewhere to go ---
but expecting is not proving, and an earlier draft asserted this in one sentence. We therefore state
the theorem for the first-order fragment and record the binding case as
Problem~\ref{prob:binding}. This is the largest gap between what the note proves and what it would like
to prove.
\end{remark}

\subsection{Idempotence}
\label{sec:idem}

\begin{proposition}[Idempotence]
\label{prop:idem}
$\bis_{\Gg^{++}}$ and $\bis_{\Gg^+}$ agree on $\Terms(\Sigma^\Sigma)$.
\end{proposition}

\begin{proof}
Since $\Gg^{++}$ has more rules than $\Gg^+$, monotonicity (Proposition~\ref{prop:mono}) gives
$\bis_{\Gg^{++}} \subseteq \bis_{\Gg^+}$. For the converse we exhibit a bisimulation rather than
invoke Theorem~\ref{thm:main} on both sides, which would be a domain mismatch: the theorem speaks of
$\Sigma$-terms under $\Gg^+$, not of $\Sigma^\Sigma$-terms under $\Gg^{++}$.

Take $\mathcal{R} = \bis_{\Gg^+}$ and check the additional transitions. The constructors adjoined by
Definition~\ref{def:ext} are of two kinds. The probe atoms are nullary: their $\mathrm{ask}$ rules have
targets $\args_{\Ask_C}()$ with no components, so the transition carries no information beyond the
identity of the atom, which $\bis_{\Gg^+}$ already fixes, since $\Kc(\Ask_C,-)$ appears in labels. The
argument formers $\args_C$ are not nullary, but the $\Get_{C,i}$ rules of $\Gg^+$ already project every
one of their components, so the second-round $\mathrm{ask}_{\args_C}$ transition is matched whenever
the first-round projections are, and adds no separating power. Hence $\bis_{\Gg^+}$ satisfies the
$\Gg^{++}$ matching conditions and is contained in $\bis_{\Gg^{++}}$.
\end{proof}

This is not a side remark. It is the reason $\Obs$ is a construction rather than the first stage of a
colimit, and it turns on a design decision that could have gone the other way: the probes are nullary
and the argument formers are exhaustively projected. Introduce a family of $n$-ary destructors with
partial projections instead and the ladder is real.

\section{Separation}
\label{sec:separation}

The operational side is done. The logical side is a separate question, and an earlier draft imported it
from the literature at a generality the cited results do not support: those are theorems about
particular spatial logics and particular calculi, and since the point of this note is to generalise
past the calculus-by-calculus treatment, borrowing the key step defeats the purpose. We prove what is
needed directly.

\begin{theorem}[Structural separation]
\label{thm:sep}
Let $E$ be admissible. Then for closed $P,Q \in \Terms(\Sigma)$,
\[
P =_{\Lo_{\mathrm{s}}(\Gg)} Q \iff P =_E Q ,
\]
where $\Lo_{\mathrm{s}}(\Gg)$ is the structural fragment.
\end{theorem}

\begin{proof}
($\Leftarrow$) Satisfaction of structural connectives is defined up to $=_E$.

($\Rightarrow$) We construct a characteristic formula. For closed $P$, by (E1) choose a lean
representative $C(t_1,\dots,t_n)$ and set
\[
\chi_P \;:=\; C(\chi_{t_1},\dots,\chi_{t_n}),
\]
which is well defined by induction on $\mu$, using (E2) for the descent, with $\chi_P := C$ for nullary
$C$.

We claim $Q \models \chi_P$ iff $Q =_E P$. If $Q \models C(\chi_{t_1},\dots,\chi_{t_n})$ then
$Q =_E C(s_1,\dots,s_n)$ with $s_i \models \chi_{t_i}$, so by induction $s_i =_E t_i$, so
$Q =_E C(\vec t) =_E P$. Conversely $P \models \chi_P$ by construction. Hence if $P \neq_E Q$ then
$\chi_P$ separates them.
\end{proof}

\begin{remark}[What this buys]
Three things. The note becomes self-contained: no separation result is imported. The image-finiteness
objection dissolves, because the argument goes from the logic to $=_E$ directly by characteristic
formulae and never passes through a Hennessy--Milner characterisation of bisimulation, which is where
image-finiteness or infinitary conjunction would have been needed. And the hypotheses are visible:
separation holds exactly for admissible $E$, so a presentation whose equational theory fails
(E1)--(E2) is one for which the note claims nothing. Note also that no boolean structure was used ---
$\chi_P$ is purely structural.
\end{remark}

\begin{corollary}[Logical correspondence]
\label{cor:adequacy}
Under the hypotheses of Theorems~\ref{thm:main} and \ref{thm:sep}, for closed $\Sigma$-terms
\[
=_{\Lo_{\mathrm{s}}(\Gg)} \;=\; =_E \;=\; \bis_{\Gg^+} .
\]
The structural fragment generated from $\Gg$ is adequate for context-labelled bisimulation in $\Gg^+$.
\end{corollary}

\section{Non-vacuity}
\label{sec:nonvacuity}

Corollary~\ref{cor:adequacy} would be worthless if $\bis_\Gg$ were already $=_E$, since the
construction would then close a gap that was not there. Remark~\ref{rem:whichbisim} also owes the
reader claim (B) for \emph{our} relation rather than the literature's. Both debts are settled by the
same observation.

\begin{proposition}[The gap is generically nonempty]
\label{prop:nonvac}
Let $\Gg$ contain two distinct nullary constructors $c,d$ of the same sort occurring in no rule of $R$
and in no equation of $E$. Then $c \bis_\Gg d$ and $c \neq_E d$, so $\bis_{\Gg^+} \subsetneq \bis_\Gg$.
\end{proposition}

\begin{proof}
Neither $c$ nor $d$ occurs in any left-hand side, so neither contributes a redex, and for any context
$D$ the transitions of $D[c]$ and $D[d]$ are in label-preserving bijection with those of $D[-]$'s own
redexes. Hence $\{(c,d)\} \cup {=_E}$ extends to a bisimulation. They are $=_E$-distinct by hypothesis,
so by Theorem~\ref{thm:main} they are not $\bis_{\Gg^+}$-related, witnessed by $\Kc(\Ask_c,-)$.
\end{proof}

The example is deliberately austere: it shows that the gap is nonempty for essentially any signature
carrying inert data, without relying on a coincidence of the equational theory. Richer witnesses are
presentation-dependent, and it is worth being explicit that they are.

\begin{example}[Replication, and a caveat]
In a presentation with replication as a primitive constructor and $E$ \emph{not} containing the
unfolding law, $!P$ and $!P \mid P$ are behaviourally indistinguishable and $=_E$-distinct, so they
witness the gap. If $E$ does contain the unfolding law they are $=_E$-equal and witness nothing. The
size of the gap is a property of the presentation, not of the calculus informally described.
\end{example}

\begin{remark}[Why minimality matters here]
Proposition~\ref{prop:nonvac} depends on labels being minimal. Were labels arbitrary enabling contexts,
a label would carry essentially the whole surrounding term, matching would force near-syntactic
agreement, and $\bis_\Gg$ would sit at or near $=_E$ --- leaving no gap to close and no content to
Corollary~\ref{cor:adequacy}. The choice of label discipline is therefore not a technical preliminary
but a load-bearing part of the statement.
\end{remark}

\section{The dial}
\label{sec:dial}

Nothing forces the construction to be run on the whole signature, and the point of
Definition~\ref{def:ext} taking $\Ob$ as a parameter is that the observer's instruments are part of the
data.

\begin{proposition}[Monotonicity]
\label{prop:mono}
If $\Ob \subseteq \Ob'$ then $\bis_{\Obs(\Gg,\Ob')} \subseteq \bis_{\Obs(\Gg,\Ob)}$, with
$\bis_{\Obs(\Gg,\varnothing)} = \bis_\Gg$ and, under the hypotheses of Theorem~\ref{thm:main},
$\bis_{\Obs(\Gg,\Sigma)} = {=_E}$ on closed $\Sigma$-terms.
\end{proposition}

\begin{proof}
The administrative rules of $\Obs(\Gg,\Ob')\setminus\Obs(\Gg,\Ob)$ carry fresh probe atoms, hence fresh
labels, and add transitions that must be matched; the largest relation satisfying more conditions is
contained in the largest relation satisfying fewer. The endpoints are Proposition~\ref{prop:still} and
Theorem~\ref{thm:main}.
\end{proof}

This is the sharp form of the resolution promised in \S\ref{sec:intro}. There is no single
observational bisimulation against which a spatial logic must be measured. Observational equivalence is
indexed by the destructuring capabilities admitted to the observer, and (A) and (B) name the two ends
of the resulting family.

What does \emph{not} follow, and what an earlier draft asserted as an immediate corollary, is that the
fragment of the logic retaining structural connectives only for $\Ob$ is adequate for
$\bis_{\Obs(\Gg,\Ob)}$.

\begin{conjecture}[Partial characterisation]
\label{conj:partial}
For $\Ob \subseteq \Sigma$, $\bis_{\Obs(\Gg,\Ob)}$ coincides with logical equivalence in the fragment
of $\Lo(\Gg)$ whose structural connectives are those of $\Ob$.
\end{conjecture}

The obstacle is real. A formula may inspect an opened constructor sitting underneath an unopened one;
equations may relate terms whose roots lie on opposite sides of $\Ob$; and a binder may move material
across an apparent structural boundary. None of that is handled by relativising
Theorem~\ref{thm:main}. We prove one endpoint-adjacent case and leave the rest.

\begin{proposition}[Downward closure]
\label{prop:downward}
Conjecture~\ref{conj:partial} holds when $\Ob$ is closed under the subterm relation of lean
representatives --- that is, when every constructor that can occur beneath a constructor of $\Ob$ in a
lean term is itself in $\Ob$.
\end{proposition}

\begin{proof}
Under that condition the descent in the proofs of Theorems~\ref{thm:main} and \ref{thm:sep} never
leaves $\Ob$, so both arguments run verbatim with $\Sigma$ replaced by $\Ob$ and $=_E$ replaced by its
restriction to the $\Ob$-generated subalgebra.
\end{proof}

We suspect Conjecture~\ref{conj:partial} is the most interesting statement in the vicinity, and
that its proof would say something about how partial structural knowledge composes. We do not give it
away as a corollary.

\section{Two cautions for downstream use}

\subsection{Transitions and derivations}
\label{sec:mult}

Lemma~\ref{lem:exposure} says the $\Ask_C$-labelled transitions of $P$ enumerate its decompositions.
The transition relation is extensional: $P \xrightarrow{c} P'$ holds when a witnessing $E$-match exists.
If two distinct $E$-matches yield the same label and the same target up to $=_E$, they are one
transition and two derivations.

For bisimulation this is immaterial --- only existence is used. For any quantitative lift it is not:
weights and costs attach to derivations, and a construction that collapses derivational multiplicity
will misprice. Anyone composing this note with a weighted or cost-accounted refinement should carry the
derivation structure explicitly rather than the relation.

\subsection{Cost is bounded by depth, not equal to it}
\label{sec:cost}

In a cost-accounted GSLT \cite{CostMonad} every rewrite carries a price, and administrative rules are
rewrites like any other. It is tempting to conclude that the price of evaluating a structural formula of
depth $d$ is a function of $d$, so that pricing schedules stipulated elsewhere become derived step
counts. That is too quick. Nesting depth bounds the length of the longest \emph{sequential}
destructuring path, but the cost of checking depends further on branching: a conjunction requires two
subchecks, a negation may require exhausting a search, and a separating conjunction may enumerate many
$E$-decompositions before finding one that works, or all of them before concluding that none does.
Sharing and caching move the figure again.

The defensible statement is that depth lower-bounds and structurally bounds sequential destructuring
depth, while actual checking cost depends additionally on branching and on the evaluation strategy for
the boolean layer. That is more useful downstream than an identity that would not survive contact with
an implementation.

\section{The construction is not capability-safe}
\label{sec:security}

We state this plainly rather than hedging it, because a reader who knows object-capability discipline
will see it in Definition~\ref{def:ext} and will want to know whether we noticed.

Unrestricted opening destroys encapsulation. Any observer may form $\Kc(\Ask_C, -)$, so any term
whatever may be taken apart, and $\Get_{C,i}$ then yields the arguments as data. In a calculus where an
opaque reference is the unit of authority, this is fatal: open the reference and harvest whatever
authority its contents carry. Confinement fails, and with it every argument that depends on it.

The repair is known in outline. Mike Stay's suggestion is that constructors and destructors should each
take a capability token, so that the right to build and the right to open are held rather than ambient.
There are two readings --- a \emph{permission}, licensing the rule, or a \emph{brand}, so that building
stamps the term and opening checks the stamp --- and the second is the more useful, since it makes
legibility a property of the term rather than a global property of the world. The token must gate the
whole kit rather than opening alone, since an observer able to take an administrative step should be
able to take its converse.

Tokens require a factory, and two observations are worth recording before we stop.

\begin{remark}[The factory is an import]
\label{rem:import}
Remark~\ref{rem:noturing} obtained the argument formers by free adjunction. That cannot be done again
here. Unforgeability is not a claim about what can be computed, nor about what can be adjoined, but
about what a context can \emph{write down}: in a signature of freely adjoined constructors every term
is writable by anybody, which is precisely why free adjunction was safe for $\args_C$ and is useless for
a token. A genuine source of privacy must be imported --- restriction, in the general case, giving a
factory of the shape $\mathsf{let}\ c = \mathsf{freshCap}()\ \mathsf{in}\ P$, essentially $(\nu c)P$.
Reflective calculi, in which a name is a quoted process, may be able to discharge the import internally;
whether quotation \emph{originates} privacy or merely propagates it from a seed we do not settle.
\end{remark}

\begin{remark}[Factories resist the construction]
\label{rem:nodestructor}
Whatever the factory is, Definition~\ref{def:ext} must not be applied to it. Give it an opening rule and
a token can be taken apart; if it can be taken apart it can be rebuilt; if it can be rebuilt it can be
forged; and if it can be forged nothing is a capability. So the capability-safe version of the
construction is not ``the observer extension, with a caveat'' but ``the observer extension on
$\Ob \subseteq \Sigma \setminus F$, where $F$ is the minting apparatus'' --- and the exclusion is what
makes the rest of it safe. An earlier draft said the factory is the \emph{unique} constructor admitting
no destructor. That is more than we can support: a system may carry several fresh-name generators,
sealed abstract types, or independent minting authorities, and in nominal presentations freshness is a
binder rather than a constructor at all. The property that matters is not uniqueness but that the
mechanism generating unforgeable authority is not itself generically destructible by the observer
extension.
\end{remark}

\subsection{Reflection, and an inherited restriction}
\label{sec:reflection}

Remark~\ref{rem:areal} recorded that congruence in the \rhoc\ calculus holds only for a restricted class
of contexts, excluding holes beneath the quote, because communication turns on name equality rather than
on bisimilarity. That restriction is inherited here and constrains $\Ob$.

The opening rule for the quote, $\Kc(\Ask_@, @(p)) \to \args_@(p)$, does not itself place a hole beneath
$@$: the hole is beside the quote, not inside it. So we do not expect it to disturb the congruence
result. But the observation is a genuine addition to the context class, and whether $\mathcal{A}$
extended by the administrative contexts still supports congruence is a question for \cite{BIHOL}'s
machinery rather than for assertion here. We flag it as Problem~\ref{prob:reflection} and note that a
reader wanting a safe default may take $\Ob$ to exclude the quote, at the cost of leaving the name
predicate outside the correspondence of Corollary~\ref{cor:adequacy}.

\section{What this settles, and for whom}
\label{sec:settles}

\paragraph{Adequacy claims survive contact with intensionality results.} A generated spatial--behavioural
logic may be asserted adequate, provided the assertion names the relation: bisimulation in the observer
extension for a stated $\Ob$, not in the base. Assertions that omit the index are not wrong so much as
ambiguous, and the ambiguity is what invites the objection.

\paragraph{Structural observation is representable as interaction.} Not identical to it. The
distinction matters: the enlargement adds transitions whose purpose is to expose structure, and an
architecture that appeals to this construction should say that its observer holds instruments, not that
its instruments were never instruments.

\paragraph{Three restrictions, kept apart.} An architecture in which an agent observes its environment
faces three different limits, and collapsing them causes trouble. There is the \emph{ideal observer}
with the full observation signature, whose reach is $=_E$. There is the \emph{capability-limited}
observer, holding some $\Ob \subsetneq \Sigma$, whose reach is the corresponding point of
Proposition~\ref{prop:mono}. And there is the \emph{budget-limited} observer, who holds instruments it
cannot afford to use, whose reach is smaller again and is not characterised here at all. Adequacy is a
claim about the first.

\paragraph{Spatial and behavioural addresses coincide.} In $\Obs(\Gg,\Ob)$ a decomposition site is a
redex position, so constructions identifying a structural locus with a modal label are well typed rather
than analogical. This does not supply an enabling theorem: a locus that is a redex position for an
administrative rule need not correspond to a \emph{proper} step reachable in $\Gg$, and claims about
reachability still require their own argument.

\paragraph{Morphisms.} Where the morphisms of a category of GSLTs preserve bisimulation \cite{Omnibus},
an assignment depending on spatial structure is not functorial, and the usual repair is to narrow the
morphism class by hand to maps preserving the generated logic. Under Theorem~\ref{thm:main} the
narrowing is not a narrowing: logic-preserving maps of $\Gg$ are exactly bisimulation-preserving maps of
$\Gg^+$, and the repair becomes functoriality on the image of $\Obs(-,\Sigma)$.

\section{Open problems}
\label{sec:open}

\begin{enumerate}
\item \label{prob:binding} \textbf{Binding.} Extend Theorem~\ref{thm:main} to binding constructors:
$\alpha$-equivalence within $E$ and the resulting measure, transitions beneath binders, and the meaning
of a probe applied to an abstraction. See Remark~\ref{rem:binding}.

\item \textbf{Conjecture~\ref{conj:partial}.} The partial characterisation, beyond the downward-closed
case of Proposition~\ref{prop:downward}.

\item \label{prob:reflection} \textbf{Reflection and congruence.} Whether the admissible context class
of \cite{BIHOL} remains a congruence class when extended by the administrative contexts, in particular
for the quote. See \S\ref{sec:reflection}.

\item \textbf{Canonicity as a universal property.} We have called the construction uniform and
presentation-directed, not canonical: it is mechanical in the signature, but $\Ob$ is a choice, the
probe indexing is a choice, and the treatment of binders will be a choice. Is $\Obs(-,\Sigma)$ a left
adjoint --- free among extensions making $\Lo_{\mathrm{s}}(\Gg)$ behavioural? An adjunction would
upgrade ``evident'' to ``forced'', which matters wherever the absence of a design choice is doing
argumentative work.

\item \textbf{The capability-indexed family.} With tokens, bisimilarity becomes indexed by the
capabilities an observer \emph{holds}, an object of the calculus rather than a parameter of the
metatheory. This is a sharper form of \S\ref{sec:dial} and needs the factory of
Remark~\ref{rem:import}. It also wants a no-amplification theorem: one cannot open one's way to a token
one was not given. That statement has a circular smell --- unforgeability rests on the gate, the gate
rests on tokens --- and needs a well-foundedness argument.

\item \textbf{Weak equivalences.} Everything here is strong. Administrative steps are natural candidates
for being made silent, and \cite{BIHOL} records that weak bisimilarity is a congruence only for reactive
contexts, so the weak theory needs its own treatment.

\item \textbf{Relation to the semantic route.} \cite{CairesVieira} obtains extensionality from partial
failures. Is there a translation carrying that model into an observer extension, so that failures and
administrative openings turn out to be two presentations of one mechanism?

\item \textbf{Precedent.} Reifying constructors as tags and supplying observation actions is elementary
enough that we expect precedents in the literature on definability by experiments, applicative and
environmental bisimulation, and observational equivalence for abstract data types. A targeted pass is
owed before any priority claim is made.
\end{enumerate}

\section*{Acknowledgements}

The construction was described to the author's satisfaction in conversation with Mike Stay, whose
capability refinement is the subject of \S\ref{sec:security}. The label discipline of
\S\ref{sec:prelim} is taken from work in preparation by C.\ B.\ Wells and Mike Stay \cite{BIHOL}. This
draft owes its shape to a referee whose objections account for most of what changed between the first
version and this one.

\begin{thebibliography}{99}

\bibitem{BIHOL} C.\ B.\ Wells and M.\ Stay. \emph{Behavior in higher-order languages.} F1R3FLY.io, in
preparation.

\bibitem{LeiferMilner} J.\ J.\ Leifer and R.\ Milner. \emph{Deriving bisimulation congruences for
reactive systems.} In Proc.\ CONCUR 2000, LNCS 1877, Springer, 2000, pp.\ 243--258.

\bibitem{Sewell} P.\ Sewell. \emph{From rewrite rules to bisimulation congruences.} Theoretical
Computer Science 274(1--2):183--230, 2002.

\bibitem{LambdaLTS} P.\ Di Gianantonio, F.\ Honsell and M.\ Lenisa. \emph{RPO, second-order contexts,
and lambda-calculus.} Logical Methods in Computer Science 5(3), 2009.

\bibitem{Bigraphs} O.\ H.\ Jensen and R.\ Milner. \emph{Bigraphs and mobile processes (revised).}
Technical Report 580, University of Cambridge Computer Laboratory, 2004.

\bibitem{SangiorgiHO} D.\ Sangiorgi. \emph{Bisimulation for higher-order process calculi.} Information
and Computation 131(2):141--178, 1996.

\bibitem{CairesObs} L.\ Caires. \emph{Behavioral and spatial observations in a logic for the
$\pi$-calculus.} In Proc.\ FOSSACS 2004, LNCS 2987, Springer, 2004.

\bibitem{CairesCardelliI} L.\ Caires and L.\ Cardelli. \emph{A spatial logic for concurrency (Part I).}
Information and Computation 186(2), 2003.

\bibitem{CairesCardelliII} L.\ Caires and L.\ Cardelli. \emph{A spatial logic for concurrency
(Part II).} Theoretical Computer Science 322(3):517--565, 2004.

\bibitem{SangiorgiExtInt} D.\ Sangiorgi. \emph{Extensionality and intensionality of the ambient
logics.} In Proc.\ POPL 2001, ACM Press, 2001.

\bibitem{HLS02} D.\ Hirschkoff, E.\ Lozes and D.\ Sangiorgi. \emph{Separability, expressiveness and
decidability in the ambient logic.} In Proc.\ LICS 2002, IEEE Computer Society, 2002.

\bibitem{HLS03} D.\ Hirschkoff, E.\ Lozes and D.\ Sangiorgi. \emph{Minimality results for the spatial
logics.} In Proc.\ FSTTCS 2003, LNCS 2914, Springer, 2003.

\bibitem{LozesAdjunct} E.\ Lozes. \emph{Adjunct elimination in the static ambient logic.} In Proc.\
EXPRESS 2003, ENTCS, Elsevier.

\bibitem{LozesThesis} E.\ Lozes. \emph{Expressivit\'e des logiques spatiales.} Th\`ese de doctorat,
ENS Lyon, 2004.

\bibitem{CairesLozes} L.\ Caires and E.\ Lozes. \emph{Elimination of quantifiers and undecidability in
spatial logics for concurrency.} In Proc.\ CONCUR 2004, LNCS 3170, Springer, 2004.

\bibitem{DGG} A.\ Dawar, P.\ Gardner and G.\ Ghelli. \emph{Adjunct elimination through games in static
ambient logic.} In Proc.\ FSTTCS 2004, LNCS 3328, Springer, 2004.

\bibitem{CairesVieira} L.\ Caires and H.\ T.\ Vieira. \emph{Extensionality of spatial observations in
distributed systems.} In Proc.\ EXPRESS 2006, ENTCS 175(3):131--149, Elsevier, 2007.

\bibitem{TuostoVieira} E.\ Tuosto and H.\ T.\ Vieira. \emph{An observational model for spatial logics.}
ENTCS 142:229--254, 2006.

\bibitem{CairesSurvey} L.\ Caires. \emph{Dynamic spatial logics: a tutorial survey.} Bulletin of the
EATCS 94, 2008.

\bibitem{Omnibus} L.\ G.\ Meredith. \emph{Graph-structured lambda theories: a reference for the working
developer.} F1R3FLY.io, 2026.

\bibitem{ScientistNote} L.\ G.\ Meredith. \emph{The mortal scientist: learning as cost-accounted
computation in a reflective higher-order calculus.} F1R3FLY.io, 2026.

\bibitem{ClassifierNote} L.\ G.\ Meredith. \emph{Classifying knot diagrams with a spatial--behavioural
logic: a logic obtained for free from the encoding of knots as processes.} F1R3FLY.io, 2026.

\bibitem{CostMonad} L.\ G.\ Meredith. \emph{Cost accounting as a monad for continued graph-structured
lambda theories.} F1R3FLY.io, 2026.

\end{thebibliography}

\end{document}
